% Generated draft: requires mathematical and source review.
\documentclass[11pt]{article}
\usepackage[margin=25mm]{geometry}
\usepackage{fontspec}
\setmainfont{Latin Modern Roman}
\usepackage{amsmath,amssymb,mathtools,unicode-math}
\setmathfont{Latin Modern Math}
\usepackage{enumitem,needspace,xurl,hyperref}
\setlength{\parindent}{0pt}
\setlength{\parskip}{5pt}
\setlength{\emergencystretch}{3em}
\newcommand{\N}{\mathbb{N}}
\newcommand{\Z}{\mathbb{Z}}
\newcommand{\Q}{\mathbb{Q}}
\newcommand{\R}{\mathbb{R}}
\newcommand{\C}{\mathbb{C}}
\newcommand{\F}{\mathbb{F}}
\newcommand{\A}{\mathbb{A}}
\newcommand{\PP}{\mathbb{P}}
\newcommand{\E}{\mathbb{E}}
\newcommand{\eps}{\varepsilon}
\newcommand{\dd}{\,\mathrm{d}}
\begin{document}
\section*{UnsolvedMath ID: 20001653}
\textbf{Source problem code:} catalogue entry
\subsection{Definitions and mathematical statement}
Let \(G=\operatorname{Out}(F_n)\) and \(\Gamma\leq G\) have finite index. Under the supplied verified reconstruction, the original shorthand ``Same questions above for subgroups of finite index'' asks to determine homomorphisms
\[
 \varphi:\Gamma\longrightarrow\operatorname{Mod}(S):
\]
can the image be infinite, especially when \(n\geq4\), must it be virtually abelian, and can the homomorphism be injective? A virtually abelian group has an abelian subgroup of finite index. The target surface and boundary conventions are not stated in the supplied short original, so use a finite-type orientable mapping class group, with boundary fixed pointwise when present, subject to checking the source's exact intended range.

The supplied correction separates extraction contamination: the heading ``Currents, Laminations, and Horofunction Boundary'' and Arnaud Hilion's Patterson--Sullivan-current normalization \(i(T,\mu_T)=1\) begin Section 5, and do not belong to Problem 4.4.

The aid gives a rank-three positive construction. For every closed orientable surface \(S\) of genus at least \(1\), there exists a finite-index subgroup \(\Gamma\leq\operatorname{Out}(F_3)\) and a homomorphism \(\Gamma\to\operatorname{Mod}(S)\) whose image is a nonabelian free subgroup. Thus, in that rank and target range, universal virtually-abelian-image and finite-image statements are false. The map is onto that free subgroup, not necessarily onto the entire mapping class group, and injectivity is not asserted.

For every \(n\geq5\), the aid reports unconditional peer-reviewed property-(T) results for \(\operatorname{Out}(F_n)\), and concludes that for every finite-index \(\Gamma\) and homomorphism \(\varphi\) to a mapping class group,
\[
 \varphi(\Gamma)\text{ virtually abelian}
 \quad\Longleftrightarrow\quad\varphi(\Gamma)\text{ finite}.
\]
Property (T) is Kazhdan's fixed-vector property: unitary representations with almost invariant vectors have nonzero invariant vectors. It passes to finite-index subgroups and quotient images in this setting. The aid further reports that any infinite image here is a nonamenable property-(T) group containing a nonabelian free subgroup \(F_2\), using the mapping-class-group subgroup alternative. This is a reduction of the virtually-abelian-image question to the finite-image question; it does not prove that all images are finite.

The same claimed equivalence for \(n=4\) is expressly conditional on the exact property-(T) theorem stated in Nitsche's computer-assisted arXiv manuscript. The source supplies neither its identifier nor the full theorem, and does not give this case the same unconditional peer-reviewed status as \(n\geq5\).

Finally, for any finite-index subgroup \(\Gamma\leq G\) of index \(m\) and homomorphism \(\varphi:\Gamma\to M\), the aid reports a coset-induced homomorphism
\[
 \widetilde\varphi:G\longrightarrow M^m\rtimes\operatorname{Sym}(m)
\]
with
\[
 \ker\widetilde\varphi
 =\operatorname{Core}_G(\ker\varphi)
 =\bigcap_{g\in G}g(\ker\varphi)g^{-1}.
\]
The construction uses coset representatives and is naturally considered up to their resulting coordinate changes. It enlarges the target to a permutational wreath product and therefore does not extend \(\varphi\) to the same target \(M\). Even an injective induced wreath-product map does not answer injectivity into the original mapping class group. The general finite-index homomorphism questions remain unresolved in the supplied account.
\subsection{Short English statement}
Ask about homomorphisms from finite-index subgroups of \(\operatorname{Out}(F_n)\) to mapping class groups. The aid constructs nonabelian free images in rank three and, unconditionally from rank five onward, equates virtually abelian images with finite images using property (T). Rank four is explicitly conditional, and coset induction changes the target.
\subsection{Sources}
\begin{itemize}
\item Local UnsolvedMath registry, permanent ID \(20001653\), code catalogue entry; canonical \texttt{aim-geometric-group-theory-notes.json}, \texttt{notes[161]}; statement status ``corrected\_verified.''
\item AIM report ``The geometry of the outer automorphism group of a free group,'' Problem 4.4, \url{https://aimath.org/WWN/outerauto/outerauto.pdf}. The supplied verification identifies the following Section 5 material as contamination; no PDF inspection is claimed here.
\item Peer-reviewed property-(T) results for ranks at least five and Nitsche's computer-assisted arXiv manuscript for rank four are described in the aid. No authors for the former, paper titles, theorem numbers, identifiers or complete authoritative references are supplied.
\item Research attempt 1: ``partial\_result,'' problem status at run ``open,'' registry status ``partially\_solved,'' rigor label ``proved,'' novelty confidence ``low.'' The machine-generated aid is not peer reviewed; independently verify before citation.
\end{itemize}
\end{document}
